% Chapter 2 — Two-body point-mass placeholder model. Follows the mandatory
% FR-29 six-section template: purpose; assumptions and domain of validity;
% derivation; discretization and implementation; validation evidence;
% references. The equation labels eq:twobody:accel and eq:twobody:firstorder
% are echoed verbatim in comments in cpp/src/models/twobody.cpp (FR-29
% equation-label-to-code traceability); renaming them breaks that trace.
\chapter{Two-Body Point-Mass Placeholder Model}
\label{ch:twobody}

\section{Purpose}
\label{sec:twobody:purpose}

This chapter documents the Phase~1 placeholder dynamics model: a spacecraft
treated as a massless test particle under the point-mass gravity of a single
central body, integrated with the classical fourth-order Runge--Kutta method.
The model exists to exercise the complete pipeline --- mission configuration,
the C++ core, the SRLOG writer, the Python loader, and the determinism gates
--- before any production physics lands, and to establish the mandatory
chapter template and equation-label-to-code traceability discipline that every
subsequent model chapter follows. It is implemented in
\texttt{cpp/src/models/twobody.cpp}.

\section{Assumptions and domain of validity}
\label{sec:twobody:domain}

Assumptions:
\begin{enumerate}
  \item The central body is a point mass or, equivalently by Newton's shell
    theorem, a spherically symmetric mass distribution; its gravitational
    field is characterized entirely by the parameter $\mu = GM$.
  \item The spacecraft is a test particle: its mass is negligible relative to
    the central body ($m \ll M$) and does not influence the motion; the
    logged mass is constant.
  \item No perturbations act: no gravity-field harmonics, no third-body
    gravity, no atmospheric drag, no solar radiation pressure, no thrust.
  \item No attitude dynamics are modeled. The logged attitude channels carry
    the fixed identity quaternion $[\,1, 0, 0, 0\,]$ (scalar-first, per
    Chapter~\ref{ch:notation}) and zero angular rate; they exist to establish
    the log schema, not to represent physics.
\end{enumerate}

Domain of validity: pipeline and infrastructure testing, determinism gating,
and Keplerian reference cases for which an analytic solution exists. The
model is \emph{not} valid for mission analysis of any real orbit regime.

Out-of-domain behavior: the model silently omits every perturbing
acceleration, so its error against a real trajectory grows secularly --- for
example, Earth-orbit nodal regression driven by the $J_2$ zonal harmonic (on
the order of degrees per day in low Earth orbit) is entirely absent. There is
no in-model detection of this condition; the guard is procedural, namely that
Phase~3 perturbation models supersede this chapter's dynamics for any
production study, and this placeholder remains selectable only as an explicit
configuration choice.

\section{Derivation}
\label{sec:twobody:derivation}

The derivation follows the standard two-body development; see
Vallado~\cite{vallado2013} for the classical treatment. Let two bodies of
masses $M$ (central body) and $m$ (spacecraft) have inertial positions
$\vv{R}_M^{I}$ and $\vv{R}_m^{I}$, and define the relative position
$\vv{r} = \vv{R}_m^{I} - \vv{R}_M^{I}$ with $r = \norm{\vv{r}}$. Newton's law
of universal gravitation gives the force on the spacecraft due to the central
body as
\begin{equation}
  \vv{F}_{m} = -\,\frac{G M m}{r^{2}}\,\frac{\vv{r}}{r},
  \label{eq:twobody:newton}
\end{equation}
and, by Newton's third law, $\vv{F}_{M} = -\vv{F}_{m}$. Applying Newton's
second law to each body in the inertial frame,
\begin{equation}
  \ddot{\vv{R}}_m^{I} = -\,\frac{G M}{r^{3}}\,\vv{r},
  \qquad
  \ddot{\vv{R}}_M^{I} = +\,\frac{G m}{r^{3}}\,\vv{r}.
  \label{eq:twobody:eachbody}
\end{equation}
Subtracting the second of equations~\eqref{eq:twobody:eachbody} from the
first yields the equation of relative motion,
\begin{equation}
  \ddot{\vv{r}} = -\,\frac{G\,(M + m)}{r^{3}}\,\vv{r}.
  \label{eq:twobody:relative}
\end{equation}
Under assumption 2 ($m \ll M$), the mass sum reduces to $M$ alone; defining
the gravitational parameter $\mu \equiv G M$ gives the model's governing
equation,
\begin{equation}
  \boxed{\;\ddot{\vv{r}} = -\,\mu\,\frac{\vv{r}}{\norm{\vv{r}}^{3}}\;}
  \label{eq:twobody:accel}
\end{equation}
with $\vv{r}$ and $\ddot{\vv{r}}$ resolved in the GCRF ($I$) frame of
Table~\ref{tab:notation:frames}.

Two conserved quantities of equation~\eqref{eq:twobody:accel} anchor the
validation evidence. First, the specific angular momentum
$\vv{h} = \vv{r} \times \vv{v}$, with $\vv{v} = \dot{\vv{r}}$, is constant:
\begin{equation}
  \dot{\vv{h}}
    = \dot{\vv{r}} \times \vv{v} + \vv{r} \times \dot{\vv{v}}
    = \vv{v} \times \vv{v}
      - \frac{\mu}{r^{3}}\,\bigl(\vv{r} \times \vv{r}\bigr)
    = \vv{0}.
  \label{eq:twobody:hconst}
\end{equation}
Second, the specific orbital energy
$\varepsilon = \tfrac{1}{2}\,\vv{v}\cdot\vv{v} - \mu/r$ is constant: using
$\dot{r} = (\vv{r} \cdot \vv{v})/r$,
\begin{equation}
  \dot{\varepsilon}
    = \vv{v} \cdot \dot{\vv{v}} + \frac{\mu}{r^{2}}\,\dot{r}
    = -\,\frac{\mu}{r^{3}}\,\bigl(\vv{v} \cdot \vv{r}\bigr)
      + \frac{\mu}{r^{3}}\,\bigl(\vv{r} \cdot \vv{v}\bigr)
    = 0.
  \label{eq:twobody:energyconst}
\end{equation}

For the circular reference case used in validation, setting the gravitational
acceleration equal to the centripetal acceleration $v_c^{2}/a$ at constant
radius $a$ gives the circular speed and period
\begin{equation}
  v_c = \sqrt{\frac{\mu}{a}},
  \qquad
  T = \frac{2 \pi a}{v_c} = 2 \pi \sqrt{\frac{a^{3}}{\mu}}.
  \label{eq:twobody:period}
\end{equation}

\section{Discretization and implementation}
\label{sec:twobody:implementation}

Equation~\eqref{eq:twobody:accel} is integrated as the autonomous first-order
system
\begin{equation}
  \vv{y} =
  \begin{bmatrix} \vv{r} \\ \vv{v} \end{bmatrix},
  \qquad
  \dot{\vv{y}} = \vv{f}(\vv{y}) =
  \begin{bmatrix} \vv{v} \\ -\,\mu\,\vv{r}/\norm{\vv{r}}^{3} \end{bmatrix},
  \label{eq:twobody:firstorder}
\end{equation}
using the shared fixed-step classical fourth-order Runge--Kutta method of
the integrator library, equation~\eqref{eq:integrators:rk4} in
Chapter~\ref{ch:integrators}, with fixed step $h = \Delta t$, local
truncation error $\mathcal{O}(h^{5})$, and global error
$\mathcal{O}(h^{4})$~\cite{hairer1993}. The model contributes only the
right-hand side of equation~\eqref{eq:twobody:firstorder}; stepping,
convergence-order evidence, and dense output are the integrator chapter's
responsibility. Because the system is autonomous, the time argument passed
through the integrator interface does not enter the stage evaluations.

Implementation notes:
\begin{enumerate}
  \item \textbf{Module and traceability.} The model lives in
    \texttt{cpp/src/models/twobody.cpp}. The source comments echo the labels
    \texttt{eq:twobody:accel} (the acceleration evaluation) and
    \texttt{eq:twobody:firstorder} (the right-hand-side assembly) verbatim,
    so each code block traces to its defining equation and the
    chapter-manifest lint can hold code and document together.
  \item \textbf{Gravitational parameter.} The Earth value is
    $\mu_\oplus = \qty{3.986004418e14}{\metre\cubed\per\second\squared}$,
    from the IERS Conventions (2010), Technical Note No.~36~\cite{iers2010}.
    The constant is defined once in the core (exposed as
    \texttt{\_core.gm("earth")}) and cited at its definition.
  \item \textbf{Determinism.} The integration loop is single-threaded with a
    fixed evaluation order, compiled with fast-math disabled and
    floating-point contraction off (decision D-10), so the same inputs on the
    same binary reproduce bit-identical trajectories --- the property gated
    by test \testid{V001}.
  \item \textbf{Logging.} States are recorded at the configured truth rate by
    decimation of integrator steps: the state is logged at $t = 0$ and at
    every $k$-th step thereafter, where $k = 1/(\Delta t \cdot
    f_{\text{truth}})$ is validated to be an exact positive integer. Logged
    values are never interpolated.
\end{enumerate}

\section{Validation evidence}
\label{sec:twobody:validation}

Table~\ref{tab:twobody:validation} lists the tests that constitute this
chapter's validation evidence. Each test identifier is machine-checked to
exist in the test tree by \texttt{scripts/lint\_chapters.py}; the tolerances
are recorded next to each test in its suite.

\begin{table}[htbp]
  \centering
  \caption{Validation evidence for the two-body placeholder model.}
  \label{tab:twobody:validation}
  \begin{tabular}{llp{6.4cm}}
    \toprule
    Test ID & Suite & Acceptance basis \\
    \midrule
    \testid{twobody_circular_orbit_analytic} &
      doctest (\texttt{cpp/tests/}) &
      Circular-orbit closure against
      equation~\eqref{eq:twobody:period}: position error after one analytic
      period $T = 2\pi\sqrt{a^{3}/\mu}$ below \qty{1}{\metre} for
      $a = \qty{6778.137}{\kilo\metre}$, $\Delta t = \qty{0.1}{\second}$. \\
    \testid{twobody_energy_momentum_drift} &
      doctest (\texttt{cpp/tests/}) &
      Relative drift of the conserved quantities $\varepsilon$
      (equation~\eqref{eq:twobody:energyconst}) and $\norm{\vv{h}}$
      (equation~\eqref{eq:twobody:hconst}) bounded over a multi-orbit
      propagation. \\
    \testid{V001} &
      \texttt{star verify} &
      Double-run determinism: two runs of the reference two-body mission
      produce SHA-256-identical SRLOG files (decision D-10). \\
    \bottomrule
  \end{tabular}
\end{table}

\section{References}
\label{sec:twobody:references}

This chapter relies on Vallado~\cite{vallado2013} for the two-body
development, on Hairer, N{\o}rsett, and Wanner~\cite{hairer1993} for the
classical Runge--Kutta method and its error orders, and on the IERS
Conventions (2010), Technical Note No.~36~\cite{iers2010} for the value of
the geocentric gravitational parameter. Full bibliographic details appear in
the bibliography.
